
由于任意流形 $V$ 都通过非紧配边 $V\times[0,1)$ 与空流形配边，除非能建立某种\BookRecordIndex{T0192}紧性，
否则定理~\ref{thm:3-1-8} 就没有用处。在 $\Sigma=S^2$ 的情形，流形
$\mathcal M(A;J)$ 本身不可能是紧的（除非它只由常值映射组成），因为 $S^2$ 的双全纯
映射所成的非紧群 $G=\operatorname{PSL}(2,\mathbb C)$ 通过重参数化作用在这个空间上。
不过，在某些情形，不带参数的 $J$-全纯球面所成的\BookRecordIndex{T0246}模空间
$\mathcal M(A;J)/G$ 是紧的。

若\BookRecordIndex{T0363}辛形式 $\omega$ \BookRecordIndex{T0372}驯服 $J$，则由\BookRecordIndex{T0249}能量恒等式（引理~\ref{lem:2-2-1}），给定同调类中
所有 \BookRecordIndex{T0003}$J$-全纯曲线的能量、从而它们的 $W^{1,2}$-范数都有一致上界。这是索伯列夫临界
情形，因此这类曲线的空间一般并不紧。事实上，标准的\BookRecordIndex{T0330}椭圆自举论证断言：若
$\mathcal M(A;J)$ 中的序列 $u^\nu$ 对某个 $p>2$ 具有一致有界的 $W^{1,p}$-范数，
则它有一个连同所有阶导数都一致收敛的子列。另一方面，如果 $u^\nu$ 的一阶导数只在
$L^2$ 中有界，那么二维能量的共形不变性便会导致\BookRecordIndex{T0244}冒泡现象；Sacks 和 Uhlenbeck 最早在
\BookRecordIndex{T0175}极小曲面的背景下发现了这一现象 \srcite{349}。确实，利用\BookRecordIndex{T0148}共形重标度的一个简单几何论证，
可以构造有限面积的 $J$-全纯映射 $v:\mathbb C\to M$；由“可去奇点”，它可延拓到
$S^2=\mathbb C\cup\{\infty\}$。这就是“球面\BookRecordIndex{T0275}气泡冒出”的现象。

下文采用的论证与 Floer 的论证非常接近。\personlater{geluomofu}{格罗莫夫}{Gromov}的
原始方法略有不同：他使用\BookRecordIndex{T0099}等周不等式和共形映射的
Schwarz 引理\sourcefootnote{原书作“Schwartz Lemma”；这里所指是共形映射的
Schwarz 引理，故据标准名称改正。}作几何论证。Pansu \srcite{320} 更详细地写出了他的
证明。下文第~\ref{sec:4-1}节和第~\ref{sec:4-5}节给出的定理~\ref{thm:4-1-2}
（可去奇点）的证明，可以让读者领略这种论证的风格。第~\ref{sec:4-3}节的\BookRecordIndex{T0266}平均值
不等式最早见于 Salamon \srcite{351}；第~\ref{sec:4-6}节和第~\ref{sec:4-7}节的证明
沿用 Hofer--Salamon \srcite{180} 对\BookRecordIndex{T0139}格罗莫夫紧性的阐述。关于这一主题的其他论述，
可见 Parker--Wolfson \srcite{321}、Ye \srcite{423} 和 Hummel \srcite{193}。近年来，
随着研究对象扩展到更复杂的全纯对象，整个主题已有显著发展。全面的处理可见
Bourgeois\sourcefootnote{原书正文拼作“Bourgeous”；同书参考文献 \srcite{46} 的作者
条目作“F. Bourgeois”，据此改正。}、Eliashberg、Hofer、Wysocki 和 Zehnder
关于\BookRecordIndex{T0352}辛场论中\BookRecordIndex{T0192}紧性结果的论文 \srcite{46}。

第~\ref{sec:4-1}节讨论 \BookRecordIndex{T0003}$J$-全纯曲线能量这一重要概念，并陈述基本的\BookRecordIndex{T0192}紧性定理和\BookRecordIndex{T0214}可去
奇点定理。第~\ref{sec:4-2}节初步介绍\BookRecordIndex{T0244}冒泡现象，并讨论若干例子；其中还证明，不带参数的
最小能量 $J$-全纯球面的\BookRecordIndex{T0246}模空间是紧的。第~\ref{sec:4-3}节证明小能量 \BookRecordIndex{T0003}$J$-全纯曲线
一阶导数的逐点估计。有了这一结果，第~\ref{sec:4-5}节中\BookRecordIndex{T0214}可去奇点定理的证明就相当容易。
不过，该证明还要用到\BookRecordIndex{T0099}等周不等式，所以第~\ref{sec:4-4}节专门讨论这个主题。
第~\ref{sec:4-6}节回到\BookRecordIndex{T0244}冒泡现象，并证明每个能量一致有界的 \BookRecordIndex{T0003}$J$-全纯曲线序列都有一个
“\BookRecordIndex{T0247}模冒泡收敛”的子列。第~\ref{sec:4-7}节给出更精细的\BookRecordIndex{T0148}共形重标度论证，由此产生一个与原
序列极限相连的\BookRecordIndex{T0275}气泡。这一结果在第 5 章\BookRecordIndex{T0340}稳定映射的\BookRecordIndex{T0139}格罗莫夫紧性证明中起关键作用。为求
完整，我们在大多数地方也讨论了圆盘的情形。

\section{能量}\label{sec:4-1}
\zhindex{nengliang}{能量}

我们先讨论\BookRecordIndex{T0249}能量恒等式。它区分了\BookRecordIndex{T0358}辛流形中的\BookRecordIndex{T0334}伪全纯曲线与一般殆复流形中的\BookRecordIndex{T0334}伪全纯曲线：
在前一种情形，能量是拓扑不变量；在后一种情形，能量没有普遍上界，因而也没有一般的
\BookRecordIndex{T0192}紧性结果。不过，我们将指出，在这一更一般的背景下仍有一些结果成立。

设 $(M,\omega)$ 是\BookRecordIndex{T0358}辛流形，配备一个 $\omega$-\BookRecordIndex{T0372}驯服\BookRecordIndex{T0088}殆复结构 $J$（见
第~\ref{sec:2-1}节）。回忆，$\omega$ 与 $J$ 在 $M$ 上确定黎曼度量
\eqref{eq:2-1-1}。定义在紧黎曼曲面
$(\Sigma,j_\Sigma,d\operatorname{vol}_\Sigma)$ 上的 \BookRecordIndex{T0003}$J$-全纯曲线
$u:\Sigma\to M$ 的能量为
\[
  E(u)
  :=\frac12\int_\Sigma |du(z)|_J^2\,d\operatorname{vol}_\Sigma
  =\int_\Sigma u^*\omega
\]
（见引理~\ref{lem:2-2-1}）。因此，定义在闭黎曼曲面上的 \BookRecordIndex{T0003}$J$-全纯曲线，其能量是
拓扑不变量。对有边界的紧黎曼曲面，只要边界映入某个\BookRecordIndex{T0219}拉格朗日子流形 $L\subset M$，
同样的结论也成立。于是，所有表示同一个同调类的 \BookRecordIndex{T0003}$J$-全纯曲线，其一阶导数都有一致的
$L^2$-上界。下文将看到，这种 $L^2$-上界不足以建立\BookRecordIndex{T0192}紧性。但是，如果一个
\BookRecordIndex{T0003}$J$-全纯曲线序列 $u_\nu$ 中各元素的一阶导数具有一致的 $L^p$-上界，那么标准的
“\BookRecordIndex{T0330}椭圆自举”技巧便可用来证明收敛子列的存在性。下面的定理给出更精确的表述；其中定义域
$\Sigma$ 可以非紧。回忆，殆复流形 $(M,J)$ 的子流形 $L$ 称为\emph{全实的}，如果
\[
  TL\cap J(TL)=\{0\}.
\]

\begin{theorem}\label{thm:4-1-1}
设 $(M,J)$ 是紧殆复流形，$L\subset M$ 是紧\BookRecordIndex{T0290}全实子流形，而 $J^\nu$ 是 $M$ 上一列在
$C^\infty$ 拓扑下收敛于 $J$ 的\BookRecordIndex{T0088}殆复结构。再设
$(\Sigma,j_\Sigma,d\operatorname{vol}_\Sigma)$ 是黎曼曲面，
$\Omega^\nu\subset\Sigma$ 是穷竭 $\Sigma$ 的递增开集列，并且
\[
  u^\nu:(\Omega^\nu,\Omega^\nu\cap\partial\Sigma)\longrightarrow(M,L)
\]
是一列 $J^\nu$-全纯曲线，满足对每个紧子集 $K\subset\Sigma$ 都有
\[
  \sup_\nu\|du^\nu\|_{L^\infty(K)}<\infty.
\]
那么 $u^\nu$ 有一个子列，在 $\Sigma$ 的紧子集上连同所有阶导数一致收敛于
\BookRecordIndex{T0003}$J$-全纯曲线
\[
  u:(\Sigma,\partial\Sigma)\longrightarrow(M,L).
\]
\end{theorem}

定理~\ref{thm:4-1-1} 的证明将在附录 B 中完成。如果映射 $u^\nu$ 对某个 $p>2$
具有一致有界的 $W^{1,p}$-范数，结论仍然成立（见定理 B.4.2）。然而，
引理~\ref{lem:2-2-1} 的\BookRecordIndex{T0249}能量恒等式只在 $p=2$ 时保证一致的 $W^{1,p}$-上界。这个
$W^{1,2}$-上界虽然不能\BookRecordIndex{T0326}推出导数一致有界，却仍对导数爆破时可能发生的情形给出一些控制。
正如第~\ref{sec:4-2}节将看到的，能量的共形不变性在此时导致\BookRecordIndex{T0275}气泡的形成。

下面陈述\BookRecordIndex{T0214}可去奇点定理；它是理解\BookRecordIndex{T0244}冒泡现象的重要分析工具。在可积情形，这是复分析中的
熟知结果。\BookRecordIndex{T0358}辛流形中 \BookRecordIndex{T0003}$J$-全纯曲线的相应证明将在第~\ref{sec:4-5}节给出。这里利用\BookRecordIndex{T0175}极小
曲面的单调性，在 $\omega$-\BookRecordIndex{T0348}相容情形给出连续性的证明。

全文用
\[
  \mathbb D:=\{z\in\mathbb C\mid |z|\le1\}
\]
表示闭单位圆盘，用
\[
  \mathbb H:=\{z\in\mathbb C\mid \operatorname{Im}z\ge0\}
\]
表示闭上半平面。

\begin{theorem}[可去奇点]\label{thm:4-1-2}
设 $(M,\omega)$ 是紧\BookRecordIndex{T0358}辛流形，$L\subset M$ 是紧\BookRecordIndex{T0219}拉格朗日子流形，$J$ 是 $M$ 上的
$\omega$-\BookRecordIndex{T0372}驯服\BookRecordIndex{T0088}殆复结构，其相伴度量为 $g_J$。
\begin{enumerate}[label=\textup{(\roman*)},leftmargin=3em]
\item 若
\[
  u:\mathbb D\setminus\{0\}\longrightarrow M
\]
是能量有限的 \BookRecordIndex{T0003}$J$-全纯曲线，即 $E(u)<\infty$，则 $u$ 可延拓为光滑映射
$\mathbb D\to M$。
\item 若
\[
  u:
  (\mathbb D\cap\mathbb H\setminus\{0\},
   \mathbb D\cap\mathbb R\setminus\{0\})
  \longrightarrow(M,L)
\]
是能量有限的 \BookRecordIndex{T0003}$J$-全纯曲线，即 $E(u)<\infty$，则 $u$ 可延拓为光滑映射
$\mathbb D\cap\mathbb H\to M$。
\end{enumerate}
\end{theorem}

\begin{proof}[第 \textup{(i)} 部分连续性的证明]
这里基本沿用\personlater{geluomofu}{格罗莫夫}{Gromov}原始工作的论证路线。我们假设
$J$ 与 $\omega$ \BookRecordIndex{T0348}相容，从而 \BookRecordIndex{T0003}$J$-全纯曲线使能量极小。如 Pansu
\srcite{320} 第 37 节所指出的，在这一情形可用\BookRecordIndex{T0175}极小曲面的\BookRecordIndex{T0092}单调性定理证明\BookRecordIndex{T0214}可去奇点定理。
该定理断言，存在常数 $c>0$ 和 $\varepsilon_0>0$（依赖于 $M$ 和度量 $g_J$），使得
对 $(M,g_J)$ 中每个边界非空的非常值紧连通\BookRecordIndex{T0175}极小曲面 $S$、每个 $x\in M$ 以及每个常数
$0<\varepsilon<\varepsilon_0$，都有（见 \srcite{232}，3.15）
\begin{equation}\label{eq:4-1-1}
  x\in S,\qquad
  \partial S\cap B(x,\varepsilon)=\varnothing
  \qquad\Longrightarrow\qquad
  \operatorname{area}_{g_J}(S)\ge c\varepsilon^2.
\end{equation}

现在令
$u:\mathbb D\setminus\{0\}\to M$ 为 \BookRecordIndex{T0003}$J$-全纯曲线\sourcefootnote{原书此处写作
“$u:\mathbb D\to M$”；但这里正在证明定理~\ref{thm:4-1-2}\textup{(i)} 中定义在
穿孔圆盘上的映射可连续延拓，且下文考察 $z\to0$ 的极限，故补回
$\setminus\{0\}$。}。对 $0<r<1$，定义闭曲线
\[
  \gamma_r:\mathbb R/2\pi\mathbb Z\longrightarrow M,
  \qquad
  \gamma_r(\theta):=u(re^{i\theta}),
\]
并用 $\ell(\gamma_r)$ 表示 $\gamma_r$ 的长度。我们断言，存在递减序列
$0<r_\nu\le1$，使得
\begin{equation}\label{eq:4-1-2}
  \lim_{\nu\to\infty}\ell(\gamma_{r_\nu})=0,
  \qquad
  \lim_{\nu\to\infty}r_\nu=0.
\end{equation}
由 $u$ 的共形性，
$|du(re^{i\theta})|=\sqrt2\,|\dot\gamma_r(\theta)|/r$。所以
\[
\begin{aligned}
  E(u)
  &=
  \int_0^1\int_0^{2\pi}
  \frac{|\dot\gamma_r(\theta)|^2}{r^2}\,
  r\,dr\wedge d\theta
  \\
  &\ge
  \int_0^1
  \left(\int_0^{2\pi}|\dot\gamma_r(\theta)|\,d\theta\right)^2
  \frac{1}{2\pi r}\,dr
  \\
  &=
  \int_0^1\frac{\ell(\gamma_r)^2}{2\pi r}\,dr.
\end{aligned}
\]
由于 $u$ 能量有限，这就证明存在满足 \eqref{eq:4-1-2} 的递减序列
$r_\nu\in(0,1]$。

由 $M$ 的\BookRecordIndex{T0192}紧性，可以取一个仍记为 $r_\nu$ 的子列和一点 $p\in M$，使得
\[
  \lim_{\nu\to\infty}\max_\theta d(\gamma_{r_\nu}(\theta),p)=0.
\]
我们断言，当 $z$ 趋于零时，$u(z)$ 收敛于 $p$。否则，存在序列 $z_k\to0$，使
$u(z_k)$ 收敛于某点 $q\ne p$。取充分小的常数 $\delta>0$，使
$\delta<\varepsilon_0$ 且 $\delta<d(p,q)$。再取充分大的 $\nu_0$，使得当
$\nu\ge\nu_0$ 时
\[
  \max_\theta d(\gamma_{r_\nu}(\theta),p)
  <\frac{d(p,q)-\delta}{2}.
\]
给定 $\nu\ge\nu_0$，取 $k$ 使
\[
  |z_k|<r_\nu,
  \qquad
  d(u(z_k),q)<\frac{d(p,q)-\delta}{2}.
\]
给定这个 $z_k$，再取 $\nu'\ge\nu_0$ 使 $r_{\nu'}<|z_k|$。记
\[
  A(r_{\nu'},r_\nu)
  :=\{z\in\mathbb C\mid r_{\nu'}\le |z|\le r_\nu\}.
\]
于是，对每个 $z\in\partial A(r_{\nu'},r_\nu)$，都有
\[
  d(u(z),u(z_k))
  \ge d(p,q)-d(p,u(z))-d(u(z_k),q)
  >\delta.
\]
在 \eqref{eq:4-1-1} 中取
$S:=u(A(r_{\nu'},r_\nu))$ 和 $x:=u(z_k)$，便得到
\[
  \int_{A(r_{\nu'},r_\nu)}|du|^2\ge c\delta^2
\]
对每个 $\nu\ge\nu_0$ 都成立。现在取递增子列
$\nu_k$（$k\ge0$），使得对所有 $k$ 都有 $\nu_k'<\nu_{k+1}$。于是\BookRecordIndex{T0170}环域
$A(r_{\nu_k'},r_{\nu_k})$（$k\ge1$）两两不交。$u$ 在每个\BookRecordIndex{T0170}环域上的能量都至少为
$c\delta^2$，这与 $u$ 能量有限矛盾。因此，在假设\BookRecordIndex{T0175}极小曲面的\BookRecordIndex{T0092}单调性定理成立的前提下，
我们已经证明 $u$ 可连续延拓到 $\mathbb D$。
\end{proof}

\begin{remark}\label{rem:4-1-3}
上述论证只在 $\omega$ 闭时成立，因为否则 \BookRecordIndex{T0003}$J$-全纯曲线不是\BookRecordIndex{T0175}极小曲面；见
引理~\ref{lem:2-2-1}。后面给出的论证更具分析性，但在证明第~\ref{sec:4-4}节的\BookRecordIndex{T0099}等周
不等式时仍使用了 $\omega$ 的闭性。Ye 在 \srcite{423} 中对任意 $\omega$ 证明了
引理~\ref{lem:4-5-1} 中局部\BookRecordIndex{T0099}等周不等式的类似结果，并借此为带埃尔米特度量的紧殆复
流形中的 \BookRecordIndex{T0003}$J$-全纯曲线建立了\BookRecordIndex{T0214}可去奇点定理。
\end{remark}

第~\ref{sec:4-2}节和第~\ref{sec:4-6}节将说明，如何结合定理~\ref{thm:4-1-1}、
定理~\ref{thm:4-1-2} 与\BookRecordIndex{T0249}能量恒等式，来理解\BookRecordIndex{T0334}伪全纯曲线的\BookRecordIndex{T0244}冒泡现象。我们将看到，取适当
子列后，\BookRecordIndex{T0243}冒泡只能发生在有限多个点附近。这是因为，非常值 $J$-全纯球面（或带拉格朗日
边界条件的 $J$-全纯圆盘）的能量不可能任意小。本节最后证明这一观察，所用的是下面
引理~\ref{lem:4-3-1} 中的先验估计。该估计对能量充分小的\BookRecordIndex{T0334}伪全纯曲线的一阶导数给出
逐点估计。

\begin{proposition}\label{prop:4-1-4}
设 $(M,J)$ 是紧殆复流形，$L\subset M$ 是紧\BookRecordIndex{T0290}全实子流形，并在 $M$ 上任取一个黎曼
度量。那么存在常数 $\hbar>0$，使得每个非常值 $J$-全纯球面
$u:S^2\to M$ 以及每个非常值 $J$-全纯圆盘
\[
  u:(\mathbb D,\partial\mathbb D)\longrightarrow(M,L)
\]
都满足
\[
  E(u)\ge\hbar.
\]
\end{proposition}

\begin{proof}
二维球面经球极投影共形微分同胚于黎曼球面
$\mathbb C\cup\{\infty\}$，二维圆盘则共形微分同胚于
$\mathbb H\cup\{\infty\}$。因此，只须对能量有限的 $J$-全纯映射
$u:\mathbb C\to M$ 和
$u:(\mathbb H,\mathbb R)\to(M,L)$ 证明结论。下面引理~\ref{lem:4-3-1} 的先验估计
断言，存在常数 $c,\delta>0$，使得对每条 \BookRecordIndex{T0003}$J$-全纯曲线 $u:\mathbb C\to M$ 和每个
$z\in\mathbb C$，都有
\[
  \int_{B_r(z)}|du|^2<\delta
  \qquad\Longrightarrow\qquad
  |du(z)|^2
  \le\frac{c}{r^2}\int_{B_r(z)}|du|^2,
\]
而对每条 \BookRecordIndex{T0003}$J$-全纯曲线
$u:(\mathbb H,\mathbb R)\to(M,L)$ 和每个 $z\in\mathbb H$，都有
\[
  \int_{B_{2r}(z)\cap\mathbb H}|du|^2<\delta
  \qquad\Longrightarrow\qquad
  \sup_{B_r(z)\cap\mathbb H}|du|^2
  \le\frac{2c}{r^2}
  \int_{B_{2r}(z)\cap\mathbb H}|du|^2.
\]
这里 $B_r(z)\subset\mathbb C$ 表示以 $z$ 为中心、半径为 $r$ 的圆盘。若取能量阈值
\sourcefootnote{原书此处写作 $E(u)<\delta$；但按本书约定
$E(u)=\frac12\int|du|^2$，而上述估计的假设是 $\int|du|^2<\delta$，故充分条件应为
$E(u)<\delta/2$。等价地也可重命名阈值；命题只需某个正下界，结论不受影响。}
$E(u)<\delta/2$，这些估计便对每个 $r>0$ 成立，故 $u$ 是常值映射。
\end{proof}

若 $\omega$ \BookRecordIndex{T0372}驯服 $J$，也可以利用可去奇点以及下一节说明的\BookRecordIndex{T0243}冒泡论证来证明
命题~\ref{prop:4-1-4}。具体地，若 $u_\nu:S^2\to M$ 是能量有界的 $J$-全纯球面序列，
则对适当选取的一列全纯映射 $\phi_\nu:\mathbb C\to S^2$，复合映射
$u_\nu\circ\phi_\nu$ 可取子列
收敛于一个能量有限的非常值 $J$-全纯平面 $u:\mathbb C\to M$。由
\BookRecordIndex{T0214}可去奇点定理~\ref{thm:4-1-2}，$u$ 延拓为 $J$-全纯球面，因而必有正能量。所以序列
$E(u_\nu)$ 不可能趋于零。对 $J$-全纯圆盘可作类似论证。不过，第~\ref{sec:4-5}节对
定理~\ref{thm:4-1-2} 的证明，使用了与上述命题~\ref{prop:4-1-4} 的证明相同的
\BookRecordIndex{T0266}平均值不等式，故这一论证归根结底并不更容易。

仍在 $\omega$-\BookRecordIndex{T0372}驯服情形，命题~\ref{prop:4-1-4} 也是下一个结果的推论；该结果断言，
给定能量上界后，只有有限多个同调类能由 $J$-全纯球面表示。利用
命题~\ref{prop:4-1-4} 和\BookRecordIndex{T0139}格罗莫夫紧性可以\BookRecordIndex{T0326}推出命题~\ref{prop:4-1-5}。不过，Joel Fish
给出了另一种远为初等的证明，现说明如下。

\begin{proposition}\label{prop:4-1-5}
设 $(M,\omega)$ 是紧\BookRecordIndex{T0358}辛流形，$J$ 是 $M$ 上的 $\omega$-\BookRecordIndex{T0372}驯服\BookRecordIndex{T0088}殆复结构。固定常数
$c>0$。那么，所有能由 $J$-全纯球面表示且满足 $\omega(A)\le c$ 的同调类
$A\in H_2(M;\mathbb Z)$ 所成的集合是有限集。
\end{proposition}

\begin{proof}
记由 $\omega$ 和 $J$ 确定的黎曼度量为
\[
  g_J
  :=\frac12\bigl(\omega(\,\cdot\,,J\,\cdot\,)
  -\omega(J\,\cdot\,,\,\cdot\,)\bigr).
\]
任取微分 $2$-形式 $\beta\in\Omega^2(M)$。对每个 $x\in M$ 和每个切向量
$\xi\in T_xM$，都有
\[
  |\beta_x(\xi,J\xi)|
  \le |\beta_x|_{g_J}|\xi|_{g_J}^2
  =|\beta_x|_{g_J}\omega(\xi,J\xi).
\]
这\BookRecordIndex{T0326}推出逐点不等式
$|u^*\beta|\le\|\beta\|_{L^\infty}u^*\omega$，从而对每条 \BookRecordIndex{T0003}$J$-全纯曲线
$u:S^2\to M$，
\begin{equation}\label{eq:4-1-3}
  \left|\int_{S^2}u^*\beta\right|
  \le
  \|\beta\|_{L^\infty}\int_{S^2}u^*\omega.
\end{equation}
现在选取有限多个闭微分 $2$-形式
$\beta_1,\ldots,\beta_\ell\in\Omega^2(M)$，使其上同调类构成
$H^2(M;\mathbb R)$ 的一组基。设 $A\in H_2(M;\mathbb Z)$ 可由 $J$-全纯球面表示，
并满足 $\omega(A)\le c$。由 \eqref{eq:4-1-3}，
\begin{equation}\label{eq:4-1-4}
  |\beta_j(A)|
  \le \|\beta_j\|_{L^\infty}c,
  \qquad
  j=1,\ldots,\ell.
\end{equation}
由于函数
\[
  H_2(M;\mathbb R)\longrightarrow\mathbb R,
  \qquad
  A\longmapsto\max_j|\beta_j(A)|
\]
是一个范数，满足 \eqref{eq:4-1-4} 的整同调类只有有限多个。
\end{proof}

下文用 $\hbar=\hbar(M,\omega,L,J)$ 表示满足命题~\ref{prop:4-1-4} 要求的最大常数
\sourcefootnote{原书下式第二个下确界的映射域写作 $(B,\partial B)$，但 $B$ 在此尚未定义，
且命题~\ref{prop:4-1-4} 明确以闭单位圆盘 $\mathbb D$ 为域；故将这个定义域符号补正为
$(\mathbb D,\partial\mathbb D)$。}：
\[
  \hbar(M,\omega,L,J)
  :=
  \min\left\{
  \inf_{\substack{u:S^2\to M\\
                  \bar\partial_J(u)=0,\ E(u)>0}}E(u),
  \quad
  \inf_{\substack{u:(\mathbb D,\partial\mathbb D)\to(M,L)\\
                  \bar\partial_J(u)=0,\ E(u)>0}}E(u)
  \right\}.
\]
若 $L=\varnothing$，则
\[
  \hbar(M,\omega,J)
  :=\hbar(M,\omega,\varnothing,J)
  =
  \inf_{\substack{u:S^2\to M\\
                  \bar\partial_J(u)=0,\ E(u)>0}}E(u).
\]
